\chapter{Materiales e Instrumental}

Para el desarrollo de este trabajo práctico se utilizó el siguiente material e instrumental:

\section{Instrumental Principal}

\begin{table}[H]
    \centering
    \begin{tabular}{|l|l|l|}
        \hline
        \textbf{Instrumento} & \textbf{Modelo} & \textbf{Especificaciones Técnicas} \\ \hline
        Osciloscopio Analógico & Pintek PS-200 & \qty{20}{MHz}, 2 canales. \\ \hline
        Osciloscopio Analógico & GW Instek GOS-623 & \qty{20}{MHz}, 2 canales. \\ \hline
        Osciloscopio Analógico & Leader 8041 & \qty{40}{MHz}, 2 canales. \\ \hline
        Osciloscopio Digital & UNI-T UTD2102CEX & \qty{100}{MHz}, \unit{1.GS/s}, 2 canales. \\ \hline
        Generador de Funciones & Instek GFG-3015 & \qty{0.01}{Hz} a \qty{15}{MHz}. \\ \hline
    \end{tabular}
    \caption{Especificaciones del instrumental utilizado.}
\end{table}

\section{Materiales Auxiliares}
\begin{itemize}
    \item Juego de puntas (dos) con atenuador X 10.
    \item Circuitos de pruebas auxiliares.
    \item Fuente de alimentación de CC (preparada específicamente para el TP).
    \item Resistor de \qty{39}{\ohm} / \qty{5}{W}.
    \item Circuito de prueba con disco telefónico.
\end{itemize}
